\documentclass[aspectratio=169]{beamer}
\input{../theme/imsb_beamer.tex}

\title{Spatial transcriptomics: foundations and the spot-based workflow}
\author{Dr.\ Robin Khatri}
\institute{Institute of Medical Systems Bioinformatics, UKE}
\date{Day 1 \textendash{} 16.06.2026}

\begin{document}

\begin{frame}[plain]\titlepage\end{frame}

\begin{frame}{Roadmap}
  \textbf{Format each day:} about 45 minutes of lecture, a short hands-on
  briefing, then two hours at the keyboard. Notebooks intentionally hold more than
  fits in two hours, so they double as a take-home reference.
  \vskip0.4em
  \begin{itemize}
    \item \textbf{Day 1.} The field, the data objects, and a full spot-based
      (Visium) workflow.
    \item \textbf{Day 2.} Imaging-based single-cell data (Xenium): segmentation,
      control metrics, and single-cell spatial statistics.
    \item \textbf{Day 3.} Tissue organisation, cell\textendash cell communication,
      and critical interpretation on a small standalone dataset.
  \end{itemize}
  \vskip0.3em
  This is a fast-moving field, so we cover the workflow \emph{and} the method
  landscape, with references to follow up.
\end{frame}

\begin{frame}{Why spatial?}
  \begin{columns}[T]
  \begin{column}{0.52\textwidth}
  Dissociation-based scRNA-seq answers \emph{which} cells are present but discards
  \emph{where} they were.
  \vskip0.5em
  Tissue function lives in organisation: who sits next to whom, gradients,
  compartments, and local circuits. Spatial transcriptomics keeps coordinates, so
  we can ask those questions directly.
  \end{column}
  \begin{column}{0.46\textwidth}
  \centering
  \begin{tikzpicture}[scale=0.9]
    \node[box,minimum width=2.6cm,minimum height=1.4cm] (t) at (0,2) {tissue\\(cells in place)};
    \node[alt,minimum width=2.6cm,minimum height=1.4cm] (d) at (0,0) {dissociated\\(location lost)};
    \node[box,minimum width=2.6cm,minimum height=1.4cm] (s) at (0,-2) {spatial assay\\(location kept)};
    \draw[arr] (t) -- node[lab,right]{scRNA-seq} (d);
    \draw[arr] (t.west) to[bend right=55] node[lab,left]{ST} (s.west);
  \end{tikzpicture}
  \end{column}
  \end{columns}
\end{frame}

\begin{frame}{The technology landscape}
  \centering
  \begin{tikzpicture}[node distance=5mm and 8mm]
    \node[box] (root) {Spatial transcriptomics};
    \node[box,below left=10mm and 6mm of root] (seq) {Sequencing-based\\(spot / barcode)};
    \node[box,below right=10mm and 6mm of root] (img) {Imaging-based\\(in situ, single molecule)};
    \node[alt,below=4mm of seq] (seql) {Visium, Visium HD\\Slide-seq, Stereo-seq};
    \node[alt,below=4mm of img] (imgl) {MERFISH, seqFISH+\\Xenium, CosMx};
    \draw[arr] (root) -- (seq);
    \draw[arr] (root) -- (img);
    \draw[arr] (seq) -- (seql);
    \draw[arr] (img) -- (imgl);
  \end{tikzpicture}
  \vskip0.4em
  Two families, two tradeoffs. We use Visium today and Xenium tomorrow as
  representatives.
  \refline{Reviews: Marx, \emph{Nat Methods} 2021 (Method of the Year 2020); Moses \& Pachter, \emph{Nat Methods} 2022.}
\end{frame}

\begin{frame}{The core tradeoff}
  \centering
  \begin{tikzpicture}[scale=1.0]
    \draw[arr] (0,0) -- (8,0) node[below right,lab]{gene throughput};
    \draw[arr] (0,0) -- (0,4.4) node[above left,lab,rotate=90,anchor=south]{spatial resolution};
    \node[box] at (6.4,0.9) {Visium\\(whole tx, spots)};
    \node[box] at (6.7,2.3) {Visium HD\\(2\,\textmu m bins)};
    \node[alt] at (2.0,3.4) {Xenium / MERFISH\\(panel, single cell)};
    \node[alt] at (3.6,1.6) {CosMx\\(panel, single cell)};
  \end{tikzpicture}
  \vskip0.3em
  Whole-transcriptome breadth and single-cell resolution still pull against each
  other. The right platform depends on the biological question.
\end{frame}

\begin{frame}{How Visium works}
  \begin{columns}[T]
  \begin{column}{0.5\textwidth}
  \begin{itemize}
    \item A slide carries a grid of barcoded spots ($\sim$55\,\textmu m). Tissue is
      placed on top, permeabilised, and mRNA is captured under each spot.
    \item A spot mixes the transcripts of the cells covering it, so a spot is a
      small \emph{neighbourhood}, not a cell.
    \item Output: a spot $\times$ gene count matrix, a stained image, and the pixel
      coordinates of each spot.
  \end{itemize}
  \vskip0.3em
  \textbf{Visium HD} replaces spots with a lawn of 2\,\textmu m squares, usually
  binned to 8\,\textmu m: near single-cell, still whole-transcriptome.
  \end{column}
  \begin{column}{0.48\textwidth}
  \centering
  \begin{tikzpicture}[scale=0.75]
    \fill[imsbblue!8,rounded corners] (-0.3,-0.3) rectangle (4.3,3.3);
    \foreach \x in {0,1,2,3,4}{\foreach \y in {0,1,2,3}{\fill[imsbblue!55] (\x,\y) circle (0.16);}}
    \draw[imsbdark,thick] (3,3) circle (0.3);
    \draw[arr] (3.3,3.0) -- (5.2,2.0);
    \node[draw=imsbdark,thick,circle,minimum size=1.6cm] (z) at (5.7,1.4) {};
    \foreach \a in {0,72,144,216,288}{\fill[imsbgrey!70] ($(5.7,1.4)+(\a:0.35)$) circle (0.12);}
    \node[lab,below=1mm of z]{one spot,\\several cells};
  \end{tikzpicture}
  \end{column}
  \end{columns}
\end{frame}

\begin{frame}{The data objects}
  \begin{columns}[T]
  \begin{column}{0.5\textwidth}
  \centering
  \begin{tikzpicture}[scale=0.9]
    \node[box,minimum width=3.6cm] (X) at (0,3) {\textbf{AnnData}};
    \node[alt,below=2mm of X,minimum width=3.6cm] (a) {X \,\,(counts)};
    \node[alt,below=1.5mm of a,minimum width=3.6cm] (b) {obs / var \,(metadata)};
    \node[alt,below=1.5mm of b,minimum width=3.6cm] (c) {obsm['spatial'] \,(coords)};
    \node[alt,below=1.5mm of c,minimum width=3.6cm] (d) {uns['spatial'] \,(image)};
    \node[alt,below=1.5mm of d,minimum width=3.6cm] (e) {obsp \,(graphs)};
  \end{tikzpicture}
  \end{column}
  \begin{column}{0.48\textwidth}
  \textbf{AnnData} holds one table: counts, per-spot and per-gene metadata,
  coordinates, image, and graphs. It is enough for one section and is what we use.
  \vskip0.5em
  \textbf{SpatialData} wraps several AnnData tables with aligned images and shapes
  in one lazy, larger-than-memory object. It is the right tool for multimodal or
  multi-sample work, and for Xenium with images.
  \refline{Wolf et al., \emph{Genome Biol} 2018 (scanpy/AnnData); Marconato et al., \emph{Nat Methods} 2025 (SpatialData).}
  \end{column}
  \end{columns}
\end{frame}

\begin{frame}{The toolkit (scverse)}
  \begin{itemize}
    \item \textbf{scanpy} / \textbf{anndata}: the single-cell core (QC, normalisation,
      clustering, plotting).
    \item \textbf{squidpy}: spatial graphs, neighbourhood statistics, image handling.
    \item \textbf{spatialdata} / \textbf{spatialdata-io}: readers and a unified object
      for images, shapes, and tables.
  \end{itemize}
  \vskip0.4em
  All interoperate, so the spatial workflow extends the familiar single-cell one.
  \refline{Palla et al., \emph{Nat Methods} 2022 (squidpy); scverse: \texttt{scverse.org}.}
\end{frame}

\begin{frame}{The analysis workflow}
  \centering
  \begin{tikzpicture}[node distance=4mm,every node/.style={box,minimum height=0.9cm,text width=1.7cm}]
    \node (qc) {QC};
    \node[right=of qc] (nm) {normalise};
    \node[right=of nm] (hv) {features (HVG)};
    \node[right=of hv] (dr) {PCA / UMAP};
    \node[below=8mm of dr] (cl) {cluster (Leiden)};
    \node[left=of cl] (an) {annotate};
    \node[left=of an] (sp) {spatial stats};
    \draw[arr] (qc)--(nm); \draw[arr] (nm)--(hv); \draw[arr] (hv)--(dr);
    \draw[arr] (dr)--(cl); \draw[arr] (cl)--(an); \draw[arr] (an)--(sp);
  \end{tikzpicture}
  \vskip0.5em
  The first six steps are the standard single-cell pipeline. The spatial layer is
  added at the end, and revisited once clusters exist.
\end{frame}

\begin{frame}{Quality control}
  \begin{itemize}
    \item Total counts and genes per spot show coverage and flag background spots.
    \item Mitochondrial fraction marks stressed or low-quality regions.
    \item Always read each metric \emph{in space}: a bad area is usually a
      contiguous patch (an edge, a fold), not scattered spots.
    \item Thresholds are dataset-specific. Look first, then choose, then check what
      you removed.
  \end{itemize}
  \figslot{a QC metric (total counts) shown on the tissue, highlighting an edge or fold artefact}{your own plot from the tutorial, or a QC figure from a methods paper}
\end{frame}

\begin{frame}{Normalisation: a landscape, not a default}
  \small
  \begin{tabular}{@{}lll@{}}
    \toprule
    Method & Idea & When \\
    \midrule
    total-count + log1p & scale to equal depth, log & fast first pass \\
    Pearson residuals & variance-stabilising GLM residuals & sharpen rare states \\
    SCTransform & regularised NB regression & batch/depth confounds \\
    \bottomrule
  \end{tabular}
  \vskip0.5em
  For spots, depth partly reflects how many cells were captured, so normalisation
  removes some density effect: keep that in mind when reading clusters.
  \refline{Hafemeister \& Satija, \emph{Genome Biol} 2019 (SCTransform); Lause et al., \emph{Genome Biol} 2021 (Pearson residuals).}
\end{frame}

\begin{frame}{Dimensionality reduction, clustering, annotation}
  \begin{itemize}
    \item PCA then a neighbour graph then UMAP, exactly as in scRNA-seq.
    \item \textbf{Leiden} clustering; resolution sets granularity, and there is no
      single true value, so choose it for the grain you need.
    \item Choose the resolution deliberately: scan a range and read a
      \emph{clustering tree} (clustree / pyclustree). Stable branches are robust
      populations; clusters that keep splitting and reshuffling signal
      over-clustering.
    \item \textbf{Annotation}: marker-based naming is transparent and reference-free;
      label transfer (celltypist, scANVI, \texttt{ingest}) scales when a matched
      reference exists.
  \end{itemize}
  \refline{Traag et al., \emph{Sci Rep} 2019 (Leiden); Dom\'inguez Conde et al., \emph{Science} 2022 (CellTypist); Xu et al., \emph{Mol Syst Biol} 2021 (scANVI).}
\end{frame}

\begin{frame}{Putting expression in space}
  \begin{itemize}
    \item A UMAP shows transcriptional groups; the spatial plot shows whether those
      groups form coherent tissue regions.
    \item A cluster that is transcriptionally tight but spatially scattered means
      something different from one that forms a compartment.
  \end{itemize}
  \figslot{clusters on the UMAP next to the same clusters on the tissue}{your own plot from the tutorial}
\end{frame}

\begin{frame}{Spatial statistics: first tools}
  \begin{columns}[T]
  \begin{column}{0.54\textwidth}
  \begin{itemize}
    \item \textbf{Spatial graph}: connect neighbouring spots; this graph underlies
      every spatial statistic.
    \item \textbf{Neighbourhood enrichment}: which clusters are neighbours more than
      chance (permutation test). Reveals tissue layering.
    \item \textbf{Spatial autocorrelation}: Moran's I (global smooth trends) and
      Geary's C (local differences) rank genes that vary across the tissue.
  \end{itemize}
  \end{column}
  \begin{column}{0.44\textwidth}
  \centering
  \begin{tikzpicture}[scale=0.8]
    \foreach \p/\x/\y/\c in {a/0/0/imsbblue,b/1.2/0.4/imsbblue,c/0.6/1.2/imsbgrey,d/2/1.4/imsbgrey,e/1.6/2.4/imsbblue,f/0.2/2.2/imsbgrey}{\fill[\c!60] (\x,\y) circle (0.18);}
    \draw[line] (0,0)--(1.2,0.4)--(0.6,1.2)--(0,0);
    \draw[line] (1.2,0.4)--(2,1.4); \draw[line] (0.6,1.2)--(1.6,2.4); \draw[line] (0.6,1.2)--(0.2,2.2);
    \node[lab,below=2mm] at (1,-0.2){spatial neighbour graph};
  \end{tikzpicture}
  \end{column}
  \end{columns}
\end{frame}

\begin{frame}{Spatially variable genes: more methods}
  \begin{itemize}
    \item Beyond Moran's I and Geary's C, dedicated models test for spatial
      structure directly: \textbf{SpatialDE} (Gaussian-process), \textbf{SPARK} and
      \textbf{SPARK-X} (scalable), and others.
    \item They differ in null model, speed, and the kind of pattern they favour
      (smooth gradient vs hotspot).
    \item For a first pass, Moran's I in squidpy is fast and interpretable.
  \end{itemize}
  \vskip0.2em
  Moran's I, the simplest test, is
  \[ I=\frac{N}{W}\,\frac{\sum_{ij} w_{ij}(x_i-\bar x)(x_j-\bar x)}{\sum_i (x_i-\bar x)^2},\quad W=\sum_{ij}w_{ij}, \]
  with $I>0$ for smooth spatial trends and null expectation $-1/(N-1)$.
  \refline{Svensson et al., \emph{Nat Methods} 2018 (SpatialDE); Sun et al., \emph{Nat Methods} 2020 (SPARK).}
\end{frame}

\begin{frame}{Spatial domains: regions, not just clusters}
  \begin{columns}[T]
  \begin{column}{0.54\textwidth}
  Expression clusters ignore \emph{where} a spot sits. A \textbf{spatial domain} is a
  spatially coherent region (a follicle, a T-cell zone) \textemdash{} coherent in both
  expression \emph{and} space.
  \vskip0.3em
  A transparent recipe (nichePCA / BANKSY): average expression over each spot's
  neighbours and stack it onto its own profile,
  \[ h_i = \big[\,x_i,\ \lambda\,\bar x_{\mathcal N(i)}\,\big], \]
  then PCA + Leiden. $\lambda$ trades fine detail for spatial coherence.
  \vskip0.3em
  Dedicated tools: \textbf{BayesSpace}, \textbf{SpaGCN}, \textbf{STAGATE},
  \textbf{GraphST}, \textbf{BANKSY}.
  \refline{Zhao et al., \emph{Nat Biotechnol} 2021 (BayesSpace); Hu et al., \emph{Nat Methods} 2021 (SpaGCN); Singhal et al., \emph{Nat Genet} 2024 (BANKSY).}
  \end{column}
  \begin{column}{0.44\textwidth}
  \centering
  \begin{tikzpicture}[scale=0.8]
    \foreach \x/\y in {0/0,0.6/0.1,0.2/0.6,0.8/0.7,1.4/0.2,1.5/0.9}{\fill[imsbblue!55] (\x,\y) circle (0.16);}
    \foreach \x/\y in {2.3/1.9,2.9/2.0,2.5/2.5,3.1/2.6,2.0/2.4,3.3/2.1}{\fill[imsbgrey!60] (\x,\y) circle (0.16);}
    \draw[line,rounded corners] (-0.3,-0.3) rectangle (1.9,1.3);
    \draw[line,rounded corners] (1.7,1.6) rectangle (3.7,2.9);
    \node[lab] at (0.8,-0.55){domain A};
    \node[lab] at (2.7,1.35){domain B};
  \end{tikzpicture}
  \vskip0.3em
  {\footnotesize Neighbour smoothing pulls scattered spots into contiguous regions.}
  \end{column}
  \end{columns}
\end{frame}

\begin{frame}{The spot problem: deconvolution}
  \begin{columns}[T]
  \begin{column}{0.54\textwidth}
  A Visium spot mixes several cells, so a ``cluster'' is a region type, not a cell
  type. To recover cell-type composition per spot, deconvolution maps a scRNA-seq
  reference onto the spots.
  \vskip0.4em
  The baseline is non-negative least squares: with a signature matrix $S$
  (genes\,$\times$\,types) and a spot $y$, solve
  $\min_{x\ge 0}\lVert Sx-y\rVert_2^2$ and normalise $x$ to proportions (we run this in
  the notebook). Production methods add a count model and platform correction.
  \vskip0.3em
  Methods: \textbf{cell2location}, \textbf{RCTD}, \textbf{Tangram},
  \textbf{stereoscope}, \textbf{SPOTlight}. Most want a GPU for realistic sizes, so
  we treat those conceptually and via precomputed results.
  \refline{Kleshchevnikov et al., \emph{Nat Biotechnol} 2022 (cell2location); Cable et al., \emph{Nat Biotechnol} 2022 (RCTD); Biancalani et al., \emph{Nat Methods} 2021 (Tangram).}
  \end{column}
  \begin{column}{0.44\textwidth}
  \centering
  \begin{tikzpicture}[scale=0.9]
    \node[draw=imsbdark,thick,circle,minimum size=1.6cm] (s) at (0,1.6) {};
    \fill[imsbblue!60] (0,1.6) ++(90:0) -- ++(0:0.7) arc (0:140:0.7) -- cycle;
    \fill[imsbgrey!60] (0,1.6) ++(0:0) -- ++(140:0.7) arc (140:250:0.7) -- cycle;
    \fill[imsbblue!30] (0,1.6) ++(0:0) -- ++(250:0.7) arc (250:360:0.7) -- cycle;
    \node[lab,above=1mm of s]{one spot = mix};
    \draw[arr] (0,0.6) -- (0,-0.2);
    \draw[imsbblue!60,line width=4pt] (-0.9,-0.6)--(-0.9,0.1);
    \draw[imsbgrey!60,line width=4pt] (-0.3,-0.6)--(-0.3,-0.1);
    \draw[imsbblue!30,line width=4pt] (0.3,-0.6)--(0.3,-0.3);
    \node[lab,below=1mm] at (-0.3,-0.7){cell-type proportions};
  \end{tikzpicture}
  \end{column}
  \end{columns}
\end{frame}

\begin{frame}{What is hard (a preview of the open problems)}
  \begin{itemize}
    \item Spot data need deconvolution; imaging data need segmentation. Both inject
      uncertainty before any biology.
    \item Targeted panels miss lowly expressed genes (cytokines, chemokines), so
      immune states are hard to call.
    \item Spatial QC metrics are still maturing; thresholds are provisional.
    \item Comparing sections means controlling for composition and batch.
  \end{itemize}
  We return to these on Days 2 and 3.
\end{frame}

\begin{frame}{Further reading}
  \footnotesize
  \begin{itemize}
    \item Moses \& Pachter. Museum of spatial transcriptomics. \emph{Nat Methods} 2022.
    \item Marx. Method of the Year 2020: spatially resolved transcriptomics. \emph{Nat Methods} 2021.
    \item Palla et al. Squidpy. \emph{Nat Methods} 2022.
    \item Marconato et al. SpatialData. \emph{Nat Methods} 2025.
    \item Squidpy and SpatialData tutorials: \texttt{squidpy.readthedocs.io}, \texttt{spatialdata.scverse.org}.
  \end{itemize}
\end{frame}

\begin{frame}{How the tutorials work}
  \begin{itemize}
    \item Two notebooks per day: \texttt{0X\_tutorial.ipynb} (worked steps plus blank
      exercises) and \texttt{0X\_solutions.ipynb} (filled).
    \item Worked cells run top to bottom, so run the whole notebook first, then do
      the exercises, which explore on copies.
    \item \emph{Method notes} in the notebook frame each choice, so it works as a
      standalone reference after the course.
  \end{itemize}
\end{frame}

\begin{frame}{Hands-on: \texttt{day\_1/01\_tutorial.ipynb}}
  Dataset: a public Visium human lymph node section. You will:
  \begin{itemize}
    \item load the section and read the data object;
    \item compute and inspect QC, then filter;
    \item normalise, select features, embed, and cluster;
    \item plot clusters and markers on the tissue;
    \item run neighbourhood enrichment and Moran's I.
  \end{itemize}
  Exercises: choose thresholds, try Pearson-residual normalisation, read a cluster
  spatially, and compare Geary's C with Moran's I.
\end{frame}

\end{document}
